\minitopic{把一条流畅回答拆成六层认识链}数字系统常把检索、筛选、生成和呈现压成一个对话框。核验时应逆向拆解：第一层是待判断的精确命题；第二层是系统实际检索到的材料；第三层是材料自身的作者、版本和证据；第四层是模型怎样选择、拼接或概括材料；第五层是界面怎样显示引用、置信和警告；第六层是使用者据此采取什么行动。上游正确不能保证下游正确，反之亦然。真实来源可能不支持回答，正确综合也可能被过度确定的界面呈现，低风险摘要又可能被带入高风险决定。

\begin{center}
\begin{tabularx}{\textwidth}{p{0.13\textwidth}p{0.27\textwidth}X}
\toprule
层级 & 典型失败 & 最低审计动作\\
\midrule
命题 & 时间、地点或主体含混 & 改写成可证伪的单句并记录截止时间\\
检索 & 漏掉关键材料或只返回同源转引 & 保存检索式，寻找反方和独立来源\\
来源 & 页面不存在、版本过旧、标题不支持正文 & 打开原件，核对作者、日期、段落和版本\\
综合 & 把可能写成已经，把相关写成因果 & 逐句标记由哪项材料支持及推理桥梁\\
界面 & 伪造确定感或隐藏系统限制 & 显示不确定、空缺、模型版本与更新时间\\
行动 & 把信息建议直接变成诊断、处罚或发布 & 指定复核者、否决权和升级条件\\
\bottomrule
\end{tabularx}
\end{center}

这张表的目的不是要求普通读者每次完成技术审计，而是按照风险分配检查。低风险的措辞建议可以抽查；涉及医疗、法律、公共安全或个人名誉的主张，则必须恢复来源和责任链。风险越高，系统越不能只靠“平均准确率很高”取得信任。

\minitopic{系统应当知道何时不回答}拒答不只是安全过滤，也是一种认识能力。系统可能因为缺少当前材料、来源彼此冲突、问题定义不清或任务超出验证能力而暂停。成熟拒答应说明是哪一类限制，并给出可执行的下一步：要求补充时间地点、列出冲突来源、建议查询权威记录或转交有权限的人。只说“我不能回答”无法帮助探究，硬给一个流畅猜测则把无知伪装成知识。

选择性回答还必须接受校准。若系统只在容易任务上回答，它的回答准确率可能很高，却把大量困难案例转给人；评价时应同时报告覆盖率、拒答后的工作量和被拒案例分布。若某些语言或群体更常被判定为“不确定”，看似谨慎的机制也可能造成不平等。总指标必须与分组失效和实际后果共同审查\parencite{buolamwini2018,nist2024}。

\minitopic{“人工在环”必须是一组可检验权限}在流程图中增加一个人工签名，并不产生有效监督。审核者至少需要看到模型输入、关键来源和不确定信息；需要有足够时间和领域能力；需要能拒绝、暂停或改变建议；还需要在拒绝自动建议后不受不合理惩罚。若人只能在数秒内点击确认，或否决会降低绩效分数，他承担了责任外观，却没有控制能力。

监督还应能改变系统。审核者发现错误后，事件是否进入记录，类似案例是否被重新检查，模型或提示版本是否更新，受影响者是否获得更正？没有反馈回路的人工审核只能拦截个别错误，不能使系统学习。有效监督因此既是个体认知任务，也是组织设计问题。

\minitopic{知识应归属于机器、团队还是机构}把系统输出称为“机器知道”，可以强调模型保存和调用信息的能力，却面临信念、理解和责任条件是否满足的问题；把知识全部归给最终使用者，又会忽略个人无法独立重建的训练、检索和工程贡献。第三种选择是把某些知识归属于人机团队或机构，但团队归属也不能仅凭成员结果相加：系统必须有稳定的信息整合、错误发现和责任分配结构。

不同任务可以给出不同归属。数据库精确返回航班时间时，机构可能拥有并传递知识；生成系统综合争议文献时，更适合说它提供候选解释，最终知识归属取决于使用者是否核验来源和理解限定。归属词不是礼貌称号，它决定谁可以作证、谁应解释理由、谁在错误后负责修正。可靠主义、能力论和社会认识论会给出不同答案，而真实制度必须把这些差异转化为权限与责任。

\minitopic{从事故复盘到可纠错制度}一次错误回答发生后，只修正输出文本不足以防止复发。事故记录应保存原问题、系统与知识库版本、检索结果、实际输出、使用情境、人工干预和造成的后果；随后区分错误属于数据、检索、推理、界面还是权限设计。若无法复现当时状态，机构就不能判断修复是否有效。来源标准和内容凭证可以帮助确认材料来自何处，但不能证明内容为真或推理正确\parencite{c2pa2025}。

制度还需要失败测试。上线前主动构造时间敏感、来源冲突、少数语言、诱导引用和高压决策案例；上线后监测新分布、保存异议并允许外部报告。测试目标不是宣布系统“通过”，而是建立一组触发降级、暂停和重新验证的条件。数字认识论的核心由此不是识别某一次幻觉，而是设计一种让错误留下痕迹、让反证能够进入、让责任不在界面中消失的认识制度。

\begin{argumentbox}
高风险人机流程的最低闭环是：命题可定位，来源可打开，模型和界面版本可追踪，人工拥有真实否决权，错误能够通知受影响者，修复结果能够用同类案例复验。少一环，系统可能仍有用，却不应被描述成已经建立了可靠的知识生产程序。
\end{argumentbox}

\index{选择性回答}
\index{人工在环}
\index{事故复盘}
